\section{Two-ended physical traces on the short component}
\label{sec:v98-paired}
For a maximal original segment write
$f_\sigma^\varepsilon(u)=A_\sigma(u)W^\varepsilon(\Psi_\sigma(u))$.
Its nonnegative one-sided traces are denoted by
$f_\sigma^-$ and $f_\sigma^+$. These are the complete physical
traces in \eqref{eq:v94-signed-input}; a critical entrance trace
need not be zero. The comparison below neither pools the occupation
classes nor identifies coincident faces from different incident
components.

Let $H$ be any Borel target restriction supported on
$0<\ell\le\delta$. Define its paired trace source at the original
roof by
\begin{equation}\label{eq:v98-paired-trace}
 T_H^\varepsilon(t)=\int
   \one_{\{a_\sigma<t<b_\sigma\}}\one_H(\Psi_\sigma(t))
                \min\{f_\sigma^-,f_\sigma^+\}\,d\sigma.
\end{equation}
The selector is evaluated at the target in both quantities being
compared. In particular it can be the old low-clock residual, which
need not be a smooth or monotone selector along a segment.

\begin{theorem}[Complete short-source and paired-trace comparison]
\label{thm:v98-paired-trace}
At every finite count and almost every roof,
\begin{equation}\label{eq:v98-paired-finite}
 |R_H^\varepsilon(t)-T_H^\varepsilon(t)|
 \le C(1+\varepsilon^{-1})\int_{t-\delta}^{t+\delta}v^0(u)\,du,
 \qquad v^0\,dt=F_\#(I_mz^{-1}\bar\nu).
\end{equation}
This untruncated positive source is finite at each count by
Theorem~\ref{thm:v98-uncut-insertion} and the bounded-gradient part.
Under \eqref{eq:v98-spectral-capacity},
\begin{equation}\label{eq:v98-paired-height}
 \limsup_m\sup_{R,n,k}m^2
       \|R_H^\varepsilon-T_H^\varepsilon\|_\infty
                         \le C(1+\varepsilon^{-1})\delta.
\end{equation}
The constants are independent of the Borel target selector. In
particular, the entire selected source on short segments with at
least one grazing endpoint satisfies the same height bound.
\end{theorem}
\begin{proof}
Absolute continuity on the closed segment gives, for every interior
$u$ where the representative is defined,
\[
 |f_\sigma^\varepsilon(u)-\min(f_\sigma^-,f_\sigma^+)|
       \le\int_{a_\sigma}^{b_\sigma}|(f_\sigma^\varepsilon)'(v)|\,dv.
\]
Both angular factors and the original guard give
\[
 |(f_\sigma^\varepsilon)'|
 \le C A_\sigma I_m/z+C\varepsilon^{-1}A_\sigma/z
 \le C(1+\varepsilon^{-1})A_\sigma I_m/z.
\]
Multiply the first inequality by the target selector. All segments
of length at most $\delta$ containing $t$ lie in
$(t-\delta,t+\delta)$. At each intermediate roof the Borel segment
bijection counts their interiors once, so the sum of these positive
integrals is bounded by the displayed integral of $v^0$.
The absolute value of the projected difference is bounded by this
sum of absolute differences. Truncate and exhaust the original
regular source, then use Tonelli. This proves
\eqref{eq:v98-paired-finite} without taking a derivative of the
Borel selector or a trace on its artificial boundary.

Under \eqref{eq:v98-spectral-capacity}, the complete first inverse
gradient has fixed-width bound $CL$, as proved after
\eqref{eq:v98-uncut-limit}. Apply it at $L=2\delta$.
If an endpoint is grazing, its original weighted trace is zero,
so the minimum in \eqref{eq:v98-paired-trace} is zero. This proves
the final assertion. An incoming critical endpoint with nonzero
trace is not included in that zero-trace argument.
\end{proof}

This comparison does not claim that a minimum commutes with
projection, or that \eqref{eq:v98-paired-trace} is the Jordan positive
part of a signed current. It controls the error of a specific positive
paired-trace comparison by the full physical derivative. The inherited
signed forcing remains formed after projection.

\begin{theorem}[A scalar reduction with distinct geometric obligations]
\label{thm:v98-endpoint}
Choose $q=1/8$, $r=c\varepsilon(B)^4\le r_{\rm col}$,
$a=\sqrt{qr}$ and $\delta=6qr$, with $c>0$ fixed and small.
Use $H=S$ from \eqref{eq:v98-disjoint-residual} in
\eqref{eq:v98-paired-trace}. If
\eqref{eq:v98-spectral-capacity} holds, then
\begin{equation}\label{eq:v98-endpoint}
 \mathcal E_M\le\mathcal H_M(T_S^{\varepsilon(B)})
 +C_MB^{-1/192}+C(qr)^{257/256}
                   +C(1+\varepsilon(B)^{-1})qr.
\end{equation}
Under that same condition, the original pointwise theorem is
equivalent along this schedule to
$\mathcal H_M(T_S^{\varepsilon(B)})\to0$.
Neither \eqref{eq:v98-spectral-capacity} nor this paired-trace
height decay is asserted established for the full Lorentz family.
\end{theorem}
\begin{proof}
The retained finite-count scalar reduction bounds $b^\varepsilon$
by $Y^\varepsilon$ plus the already paid clock and protected
compression terms. Split $Y=C\sqcup S$ before projection.
Theorem~\ref{thm:v98-critical-certificate} pays $C$, and
Theorem~\ref{thm:v98-paired-trace} compares $S$ to $T_S$.
The retained positive remainder proves \eqref{eq:v98-endpoint}.
Every scale and auxiliary band is fixed before the collision limit;
\eqref{eq:v98-spectral-capacity} is invoked only at each such fixed
band. The band sequence is enlarged afterwards, and $B$ is last.

For the schedule, the additional powers are
$(qr)^{257/256}=O(\varepsilon^{257/64})$ and
$(1+\varepsilon^{-1})qr=O(\varepsilon^3)$; both vanish.
Sufficiency follows. Conversely $S^\varepsilon\le b^\varepsilon$,
so scalar closure and the positive remainder imply
$\mathcal H_M(S^{\varepsilon(B)})=O(B^{-1/192})$.
Equation~\eqref{eq:v98-paired-height} then implies decay of $T_S$.
This necessity statement is conditional on the same spectral-capacity
estimate as its sufficiency statement.
\end{proof}

Without \eqref{eq:v98-spectral-capacity}, the finite-count comparisons
\eqref{eq:v98-critical-certificate} and \eqref{eq:v98-paired-finite}
remain valid with the explicit error in
\eqref{eq:v98-uncut-insertion}. Thus the new reduction identifies two
separate tasks: control a weighted two-dimensional momentum-fiber
multiplicity against the spectral error, and control the original
paired positive boundary traces on genuinely short segments. The
first is not a roof-height assumption; the second retains the exact
target selector, physical endpoints and original guard. Finite-count
finiteness does not prove either asymptotic task. Scalar upper closure,
not a finite-norm consequence, still precedes every unrestricted
point-conditioned downstream assertion.
